\section{Related Work: Evaluation and Resource Sharing}
\subsection{Measurement Studies}
Breaking the Ice divides vLLM startup into six stages, measures how their costs
change with model and system parameters, and builds a predictor
\citep{breakingice}. \citet{specillusion} test speculative decoding in vLLM
across workloads, models and batch settings. They relate measured execution
times to acceptance behavior and speedup bounds. These studies support using
phase measurements and strong runtime controls. Empirical evaluation of
autoregressive inference on resource-constrained Raspberry Pi 5 nodes also
provides context for the platform used here \citep{daian2026rpi}.
We apply this approach to
ingestion and inference, check answers and retrieved versions, and fix timing
predictions before testing other cases. Our study covers one 14B model and
controlled documents. Its predictor does not show an accuracy advantage.

\emph{Beyond the Buzz} explores how workloads and configurations affect
prefill/decode disaggregation \citep{beyondbuzz}. ProfInfer uses eBPF probes to
inspect inference in engines such as llama.cpp \citep{profinfer}. These Industry
Track studies provide context for our design and measurements. ProfInfer also
shows why operator and phase behavior matter. Our process and thread samples
and server timers give less detail. They do not separate cache from bandwidth
causes. We have not run either system as a competitor in our experiments.

\subsection{Scheduling and CPU Allocation}
Sarathi-Serve uses chunked prefills and stall-free batching to reduce
prefill/decode interference. DistServe places the two phases on separate
resources and tunes allocation for their latency requirements
\citep{sarathi,distserve}. SuperInfer combines scheduling with KV transfers on
Superchips to meet response time targets \citep{superinfer}. These systems
already account for phase demand. We study the effect of a fixed CPU reservation
that remains after ingestion ends. We do not test a scheduler that returns
idle cores to inference or show that such a scheduler would be new.

PagedAttention reduces KV fragmentation and supports sharing. FlexGen places
work across GPU, CPU and disk to improve throughput with limited GPU memory
\citep{pagedattention,flexgen}. Their objectives differ from the time to finish
one interactive request. ServerlessLLM combines checkpoint loading, migration
and scheduling based on startup cost \citep{serverlessllm}. This motivates
reporting readiness separately from warm requests. Our execution baseline is
the frozen ARM llama.cpp server \citep{llamacpp}. Published GPU speedups do
not estimate what these systems would achieve within our CPU budget.

\subsection{Freshness and Answer Quality in RAG}
METIS adapts RAG configuration and scheduling to meet quality requirements
\citep{metis}. FreshDiskANN maintains a graph index under streaming updates
\citep{freshdiskann}. We study when a separately arriving question can use a
published update while inference shares the same resources. Our MiniLM worker
uses an exact dense index. We do not propose an ANN method or compare against
implementations of METIS or FreshDiskANN.

ALCE evaluates answer correctness and citation quality separately. SQuAD2
includes questions that cannot be answered from the supplied text
\citep{alce,squad2}. These works motivate our separate checks for answers,
citations and versions. The earlier 24-article test includes unanswerable
questions. E2/E3 use six answerable questions under different configurations.
An exact-match failure alone does not show a hallucination. Human review is
still needed to assess meaning and support.

\subsection{Earlier Mechanism Controls}
GPTQ and AWQ are quantization methods \citep{gptq,awq}. DroidSpeak and ElastKV
reuse KV across model variants or precisions and selectively recompute state
\citep{droidspeak,elastkv}. Attention omission has also been studied
\citep{attentiondrop}. EdgeShard addresses placement across devices, while
PipeLive supports live pipeline reconfiguration \citep{edgeshard,pipelive}.
These works provide context for our earlier precision, state and execution
controls. The E1--E3 evaluation instead centers on resource allocation.

Our contribution is the evaluation framework and the measured relationship between CPU
sharing, phase costs, fixed reservations and answer timing. All concurrent
cases remain limited by the inference slot, and the predictor performs about
as well as a simple baseline added later. The findings identify problems to
test further. They do not establish a first general evaluation method,
a new scheduler or superiority over the systems discussed here.
